\section{Гомоморфизмы алгебраических расширений}

Теперь мы перенесём \rf{2.1} и \rf{2.4} с простых расширений на расширения, порождённые несколькими корнями одного
многочлена, и выведем единственность поля разложения.

\begin{proposition}[2.12]
Пусть $f\in F[X]$. Пусть $E$~--- расширение $F$, порождённое корнями $f$ в $E$, а $\Om$~--- расширение $F$, расщепляющее
$f$. Тогда существует $F$-гомоморфизм $\varphi\colon E\to\Om$; число таких гомоморфизмов не больше $[E:F]$ и равно $[E:F]$,
если $f$ не имеет кратных корней в $\Om$.
\end{proposition}

Уточним формулировку.
\begin{itemize}[leftmargin=*]
\item <<$E$ порождено корнями $f$>> означает, что $E=F[\alpha_1,\dots,\alpha_m]$, где $\alpha_i\in E$~--- \emph{некоторые}
корни $f$, не обязательно все (можно считать их различными, так что $m\le\deg f$); $E$ не обязано расщеплять $f$. Например, годится поле корня $E=F[\alpha]$ любого
неприводимого делителя $f$. В частности, $E/F$ конечно (\rf{1.30}).
\item $\Om$~--- \emph{любое} поле, в котором $f$ раскладывается на линейные множители; оно не обязано быть порождено
корнями $f$ и может быть сколь угодно большим (например, $\CC$ при $F=\Q$).
\item Условие отсутствия кратных корней \emph{достаточно} для равенства, но не необходимо (пример \rf{2.12b}(2)).
\end{itemize}

Удобно доказать сразу более общее утверждение~--- о продолжениях гомоморфизма $\varphi_0\colon F\to\Om$, как в
\rf{2.4}; предложение \rf{2.12}~--- его частный случай, когда $\varphi_0$~--- включение $F\subset\Om$ (тогда $\varphi_0f=f$).

\begin{proposition}[2.12a]
Пусть $f\in F[X]$, $E=F[\alpha_1,\dots,\alpha_m]$, где $\alpha_i\in E$~--- корни $f$, и $\varphi_0\colon F\to\Om$~--- гомоморфизм в
поле $\Om$, в котором $\varphi_0f$ раскладывается на линейные множители. Тогда $\varphi_0$ продолжается до гомоморфизма
$E\to\Om$; число продолжений не больше $[E:F]$ и равно $[E:F]$, если $\varphi_0f$ не имеет кратных корней в $\Om$.
\end{proposition}

\begin{proof}
Можно считать $f$ унитарным (умножение на константу не меняет корней). По условию
$\varphi_0f=\prod_{j=1}^n(X-\beta_j)$ в $\Om[X]$.

\emph{Исходное наблюдение.} Пусть $L$~--- подполе $\Om$, содержащее $\varphi_0(F)$, и $g\in L[X]$~--- унитарный делитель
$\varphi_0f$ в $L[X]$. Тогда $g$ делит $\varphi_0f$ и в $\Om[X]$, и по однозначности разложения в $\Om[X]$ это произведение
некоторых из $X-\beta_j$. Значит, $g$ раскладывается в $\Om$, и если $\varphi_0f$ не имеет кратных корней, то их нет и у $g$.

\emph{Индукция по $m$.} При $m=0$ имеем $E=F$, единственное <<продолжение>>~--- сам $\varphi_0$, и $[E:F]=1$. Пусть $m\ge1$,
$E'=F[\alpha_1,\dots,\alpha_{m-1}]$, так что $E=E'[\alpha_m]$. По предположению индукции $\varphi_0$ имеет $N'\ge1$
продолжений $\psi\colon E'\to\Om$, причём $N'\le[E':F]$, и $N'=[E':F]$, если у $\varphi_0f$ нет кратных корней.

Пусть $g\in E'[X]$~--- минимальный многочлен $\alpha_m$ над $E'$. Так как $f(\alpha_m)=0$ и $f\in F[X]\subset E'[X]$, имеем
$g\mid f$ в $E'[X]$. Применяя к коэффициентам продолжение $\psi$ (оно совпадает с $\varphi_0$ на $F$), получаем
\[
\psi g\mid\psi f=\varphi_0f\quad\text{в }\psi(E')[X].
\]
По исходному наблюдению (с $L=\psi(E')$) многочлен $\psi g$ раскладывается в $\Om$; обозначим через $r_\psi$ число его
различных корней в $\Om$. Тогда $1\le r_\psi\le\deg g=[E:E']$, и $r_\psi=[E:E']$, если у $\varphi_0f$ нет кратных корней. По
\rf{2.4}(b), применённому к полю $E'$, элементу $\alpha_m$ и гомоморфизму $\psi$, у $\psi$ ровно $r_\psi$ продолжений на $E$.

Каждое продолжение $\varphi_0$ на $E$ ограничивается на $E'$ до некоторого продолжения $\psi$. Поэтому число продолжений
$\varphi_0$ на $E$ равно
\[
\sum_{\psi}r_\psi\ \le\ N'\,[E:E']\ \le\ [E':F]\,[E:E']=[E:F]
\]
(\rf{1.20}); оно не меньше $1$, а если у $\varphi_0f$ нет кратных корней, то в обоих неравенствах достигается равенство.
\end{proof}

\par\noindent\begin{minipage}{\linewidth}
\begin{center}
\begin{tikzcd}[column sep=normal,row sep=large]
F\arrow[r,hook]\arrow[drr,cold,"\varphi_0"',bend right=12] & F[\alpha_1]\arrow[r,hook]\arrow[dr,cnew,"\varphi_1" description] &
F[\alpha_1,\alpha_2]\arrow[r,hook]\arrow[d,cnew,"\varphi_2" description] & \cdots\arrow[r,hook] &
F[\alpha_1,\dots,\alpha_m]=E\arrow[dll,cnew,"\varphi_m=\varphi",bend left=12]\\
& & \Om & &
\end{tikzcd}
\end{center}
\figcap{Доказательство \rf{2.12a}: гомоморфизм продолжается по одному корню. На $i$-м шаге продолжений $\varphi_{i-1}$
столько, сколько различных корней в $\Om$ у многочлена $\varphi_{i-1}g_i$, где $g_i$~--- минимальный многочлен $\alpha_i$ над
$F[\alpha_1,\dots,\alpha_{i-1}]$; их не больше $\deg g_i$.}
\end{minipage}\par\medskip

Подчеркнём, что на $i$-м шаге в $\Om$ раскладывается не сам минимальный многочлен $g_i$ (его коэффициенты лежат в
$F[\alpha_1,\dots,\alpha_{i-1}]$, а не в $\Om$), а его образ $\varphi_{i-1}g_i$, делящий $\varphi_{i-1}f=\varphi_0f$.

\begin{example}[2.12b]
(1) $E$ не обязано быть полем разложения: $f=X^3-2$, $E=\Q[\sqrt[3]2]$, $\Om=\CC$. Гомоморфизмов $E\to\CC$ ровно
$3=[E:\Q]$ (\rf{2.2}).

(2) Равенство возможно и при кратных корнях: $f=(X^2-2)^2$, $E=\Q[\sqrt2]$, $\Om=\CC$. У $f$ кратные корни, но
$\Q$-гомоморфизмов $E\to\CC$ два ($\sqrt2\mapsto\pm\sqrt2$), т.\,е.\ $[E:\Q]$. Важно лишь, чтобы на каждом шаге многочлен
$\varphi_{i-1}g_i$ имел различные корни.

(3) Неравенство может быть строгим: $F=\F_p(T)$, $f=X^p-T$, $E=F[\alpha]$~--- его поле корня, $\Om=E$. Гомоморфизм
$E\to\Om$ один, а $[E:F]=p$ (\rf{2.3}).

(4) Без условия, что $\Om$ расщепляет $f$, гомоморфизма может не быть ($E=\Q[i]$, $\Om=\R$, $f=X^2+1$), хотя может и
быть ($E=\Q[\sqrt[3]2]$, $\Om=\R$, $f=X^3-2$: один гомоморфизм).
\end{example}

\begin{corollary}[2.13]
Если $E_1$ и $E_2$~--- поля разложения многочлена $f\in F[X]$, то каждый $F$-гомоморфизм $E_1\to E_2$~--- изоморфизм. В
частности, любые два поля разложения $f$ $F$-изоморфны.
\end{corollary}

\begin{proof}
Каждый $F$-гомоморфизм $E_1\to E_2$ инъективен и $F$-линеен, поэтому если такой гомоморфизм существует, то
$[E_1:F]\le[E_2:F]$ (оба расширения конечны, \rf{1.30}). Поле $E_1$ порождено корнями $f$, а $E_2$ расщепляет $f$, так что
по \rf{2.12} $F$-гомоморфизм $E_1\to E_2$ существует; по симметрии существует и $E_2\to E_1$. Значит,
$[E_1:F]=[E_2:F]$, и каждый $F$-гомоморфизм $E_1\to E_2$~--- инъективное линейное отображение пространств одинаковой
конечной размерности, т.\,е.\ изоморфизм.
\end{proof}

Можно рассуждать и так: если $\varphi\colon E_1\to E_2$ и $\psi\colon E_2\to E_1$~--- $F$-гомоморфизмы, то $\psi\circ\varphi$~---
инъективный $F$-линейный эндоморфизм конечномерного пространства $E_1$, а значит, биекция (размерность ядра плюс
размерность образа равна $[E_1:F]$); так же $\varphi\circ\psi$~--- биекция $E_2$, откуда $\varphi$ биективно. Ещё одно
доказательство даёт следующее утверждение (при $F'=F$, $\varphi_0=\id$): образ $\varphi(E_1)$~--- поле, порождённое над $F$
\emph{всеми} корнями $f$ в $E_2$, т.\,е.\ всё $E_2$.

\begin{corollary}[2.13a]
Пусть $\varphi_0\colon F\to F'$~--- изоморфизм полей, $f\in F[X]$, $E$~--- поле разложения $f$ над $F$, а $E'$~--- поле
разложения $\varphi_0f$ над $F'$. Тогда $\varphi_0$ продолжается до изоморфизма $E\to E'$, и число таких продолжений не
больше $[E:F]$ (и равно $[E:F]$, если у $f$ нет кратных корней).
\end{corollary}

\begin{proof}
По \rf{2.12a} (с $\Om=E'$) существует продолжение $\varphi\colon E\to E'$, и их не больше $[E:F]$; если у $f$ нет кратных
корней, то их нет и у $\varphi_0f$ (применяем $\varphi$ к разложению $f$ на различные линейные множители в $E[X]$), и
продолжений ровно $[E:F]$. Пусть $f=c\prod_i(X-\alpha_i)$ в $E[X]$. Применяя $\varphi$ к коэффициентам, получаем
$\varphi_0f=\varphi_0(c)\prod_i(X-\varphi(\alpha_i))$ в $E'[X]$, т.\,е.\ $\varphi(\alpha_i)$~--- все корни $\varphi_0f$ в $E'$. Значит,
$\varphi(E)=\varphi(F)[\varphi(\alpha_1),\dots]=F'[\text{корни }\varphi_0f]=E'$, и $\varphi$~--- изоморфизм.
\end{proof}

Итак, \emph{поле разложения многочлена единственно с точностью до изоморфизма}~--- но, как мы увидим в \rf{2.16},
выделенного изоморфизма между двумя полями разложения, вообще говоря, нет.

\begin{corollary}[2.14]
Пусть $E$ и $L$~--- расширения $F$, причём $E$ конечно над $F$. Число $F$-гомоморфизмов $E\to L$ не больше $[E:F]$.
\end{corollary}

\begin{proof}
По \rf{1.30} $E=F[\alpha_1,\dots,\alpha_m]$ с алгебраическими $\alpha_i$. Пусть $f\in F[X]$~--- произведение минимальных
многочленов элементов $\alpha_i$ над $F$; тогда $E$ порождено над $F$ корнями $f$. Будем рассматривать $f$ как многочлен
из $L[X]$ и возьмём его поле разложения $\Om$ \emph{над полем $L$} (оно существует по \rf{2.7}, применённому к полю $L$).
По построению $\Om\supset L$; обозначим через $\iota\colon L\hookrightarrow\Om$ включение. Поле $\Om$ расщепляет $f$, а $E$
порождено корнями $f$, поэтому по \rf{2.12} число $F$-гомоморфизмов $E\to\Om$ не больше $[E:F]$. Каждый $F$-гомоморфизм
$\sigma\colon E\to L$ даёт $F$-гомоморфизм $\iota\circ\sigma\colon E\to\Om$, и различные $\sigma$ дают различные
$\iota\circ\sigma$. Следовательно,
\[
\bigl|\{F\text{-гомоморфизмы }E\to L\}\bigr|\ \le\ \bigl|\{F\text{-гомоморфизмы }E\to\Om\}\bigr|\ \le\ [E:F].\qedhere
\]
\end{proof}

\par\noindent\begin{minipage}{\linewidth}
\begin{center}
\begin{tikzcd}[column sep=large,row sep=large]
E=F[\alpha_1,\dots,\alpha_m]\arrow[r,cold,"\sigma"]\arrow[rr,cnew,bend right=22,"\iota\circ\sigma"'] & L\arrow[r,cnew,hook,"\iota"] & \Om\\[-4pt]
F\arrow[u,hook] & & {\scriptstyle\text{поле разложения $f$ над $L$}}
\end{tikzcd}
\end{center}
\figcap{К доказательству \rf{2.14}: поле $\Om$ строится над $L$, поэтому $L\subset\Om$ по построению, и каждый
$F$-гомоморфизм $E\to L$ становится $F$-гомоморфизмом в поле $\Om$, к которому применимо \rf{2.12}.}
\end{minipage}\par\medskip

Здесь $\Om$~--- поле разложения $f$ именно над $L$, а не над $F$: например, при $F=\Q$, $L=\R$, $E=\Q[\sqrt[3]2]$ имеем
$f=X^3-2$ и $\Om=\CC$ (\rf{2.7a}). Само поле $L$ может не расщеплять $f$, и гомоморфизмов $E\to L$ может не быть вовсе
($E=\Q[i]$, $L=\R$). Поле $L$ произвольно, в частности, оно не обязано быть алгебраическим над $F$.

\begin{remark}[2.15]
Пусть $E_1,E_2,\dots,E_m$~--- конечные расширения $F$, а $L$~--- расширение $F$. Тогда существует конечное расширение
$\Om/L$, содержащее изоморфный образ каждого $E_i$ (т.\,е.\ для каждого $i$ есть $F$-гомоморфизм $E_i\to\Om$).

Действительно, пусть $f_1\in F[X]$~--- произведение минимальных многочленов образующих $E_1$ над $F$ (как в
доказательстве \rf{2.14}), и $L_1$~--- поле разложения $f_1$ над $L$. Расширение $L_1/L$ конечно: по \rf{2.7}
$[L_1:L]\le(\deg f_1)!$. По \rf{2.12} существует $F$-гомоморфизм $E_1\to L_1$. Затем то же рассуждение для $E_2$ и поля
$L_1$ даёт конечное расширение $L_2/L_1$ с $F$-гомоморфизмом $E_2\to L_2$, и т.\,д. Поле $\Om=L_m$ конечно над $L$ по
\rf{1.20}, и в него вкладываются все $E_i$ (через $E_i\to L_i\subset L_m$).
\end{remark}

\par\noindent\begin{minipage}{\linewidth}
\begin{center}
\begin{tikzcd}[column sep=normal,row sep=large]
F\arrow[r,hook]\arrow[d,hook] & E_1\arrow[d,cnew] & & F\arrow[r,hook]\arrow[d,hook] & E_2\arrow[d,cnew] & & & E_m\arrow[d,cnew]\\
L\arrow[r,cnew,hook,"\text{кон.}"'] & L_1 & \Rightarrow & L_1\arrow[r,cnew,hook,"\text{кон.}"'] & L_2 & \Rightarrow\ \cdots & L_{m-1}\arrow[r,cnew,hook,"\text{кон.}"'] & L_m=\Om
\end{tikzcd}
\end{center}
\figcap{Замечание \rf{2.15}: на каждом шаге $L_i$~--- поле разложения над $L_{i-1}$ произведения минимальных многочленов
образующих $E_i$; оно конечно над $L_{i-1}$ по \rf{2.7} и содержит образ $E_i$ по \rf{2.12}.}
\end{minipage}\par\medskip

\begin{aside}[2.15a]
Квадраты на этом рисунке напоминают \emph{кодекартовы квадраты} (корасслоённые произведения), но ими не являются:
поле $L_1$ и гомоморфизм $E_1\to L_1$ не единственны и не обладают никаким универсальным свойством. Пусть, например,
$F=\Q$, $E_1=L=\Q[\alpha]$, $\alpha^3=2$. По \rf{2.4c}(2) в $L[X]$ имеем $X^3-2=(X-\alpha)(X^2+\alpha X+\alpha^2)$, и
<<минимальные>> поля, в которые вкладываются одновременно $L$ и $E_1$, таковы: $L[X]/(X-\alpha)=L$ (образ $E_1$
совпадает с $L$) и $L[X]/(X^2+\alpha X+\alpha^2)$~--- поле разложения $X^3-2$ степени $6$ (образ $E_1$ отличен от $L$). Это поля
разных степеней. Роль <<универсального>> объекта здесь играет кольцо $L[X]/(X^3-2)$ (тензорное произведение
$E_1\otimes_FL$), а оно не поле: по китайской теореме об остатках
\[
L[X]/(X^3-2)\ \cong\ L\times L[X]/(X^2+\alpha X+\alpha^2).
\]
Вообще, пусть $E_1=F[X]/(g)$ и $g=\prod_{j=1}^rg_j^{e_j}$~--- разложение $g$ на попарно различные неприводимые
множители над $L$ (если $g$ сепарабелен, то все $e_j=1$). По китайской теореме об остатках
$L[X]/(g)\cong\prod_jL[X]/(g_j^{e_j})$, и каждое поле $L[X]/(g_j)$~--- одно из возможных <<$L_1$>> (при $e_j>1$
сомножитель $L[X]/(g_j^{e_j})$ не поле, но отображается на поле $L[X]/(g_j)$). Например, при $F=\F_p(T)$, $g=X^p-T$ и
$L=E_1=F[\alpha]$ имеем $g=(X-\alpha)^p$ над $L$, и возможность единственна: $L_1=L$. Конечность $L_1/L$ при этом
сохраняется при любом выборе: $[L[X]/(g_j):L]=\deg g_j\le\deg g=[E_1:F]$.
\end{aside}

\begin{remark}[2.16: предупреждение]
Пусть $f\in F[X]$. Если $E$ и $E'$~--- два поля разложения $f$, то $F$-изоморфизм $E\to E'$ существует (\rf{2.13}), но
выделенного (<<канонического>>) изоморфизма, как правило, нет, и отождествление таких полей может привести к ошибкам
и путанице. Например, $\R[X]/(X^2+1)$ и $\CC$~--- поля разложения $X^2+1$ над $\R$, и есть ровно два $\R$-изоморфизма
между ними ($x\mapsto i$ и $x\mapsto-i$, где $x$~--- класс $X$), ни один из которых ничем не лучше другого; даже при $E=E'$ имеется
нетривиальный автоморфизм (комплексное сопряжение). Поэтому, работая с абстрактным полем разложения $X^3-2$, нельзя
говорить о <<вещественном корне $\sqrt[3]2$>>: это понятие появляется лишь после выбора вложения в $\CC$.

Кроме того, бессмысленно говорить о <<поле $F[\alpha]$, порождённом корнем $f$>>, если $f$ приводим: поля,
порождённые корнями разных неприводимых множителей, никак не связаны. Так, для $f=(X^2-2)(X^2-3)$ над $\Q$ получаются
$\Q[\sqrt2]$ и $\Q[\sqrt3]$, и они не изоморфны: $\Q$-изоморфизм перевёл бы $\sqrt2$ в квадратный корень из $2$ в
$\Q[\sqrt3]$, а его там нет, так как $1,\sqrt2,\sqrt3,\sqrt6$ линейно независимы над $\Q$ (\rf{1.32a}).

Даже для неприводимого $f$ бессмысленно говорить о <<поле $F[\alpha,\beta]$, порождённом двумя корнями $\alpha,\beta$>>:
расширения $F[\alpha]$, порождённые корнями разных неприводимых множителей $f$ в $F[\alpha][X]$, могут быть совсем
разными (\rf{2.4b}--\rf{2.4d}: для $X^4-2$ над $\Q$ получаются поля степеней $4$ и $8$).
\end{remark}
